\documentclass[11pt,a4paper]{article}

\usepackage[margin=1in]{geometry}
\usepackage[T1]{fontenc}
\usepackage{lmodern}
\usepackage{microtype}
\usepackage{booktabs}
\usepackage{array}
\usepackage{amsmath}
\usepackage{amssymb}
\usepackage{hyperref}
\usepackage{enumitem}
\usepackage{xcolor}

\hypersetup{colorlinks=true, linkcolor=blue!50!black,
            urlcolor=blue!60!black, citecolor=blue!50!black}

\newcolumntype{L}[1]{>{\raggedright\arraybackslash}p{#1}}
\newcolumntype{C}[1]{>{\centering\arraybackslash}p{#1}}

\title{Scaling Up the Experiment: Before and After\\
\large What changed when we moved from a tiny test set to a large one}

\date{}

\begin{document}

\maketitle
\vspace{-2em}

\noindent
To scale the experiment up to the full dataset, I split the work into three
separate cloud jobs (I call them Kernel~A, B, and C). This section explains what
each one does and its current status. The rest of the report uses the results
from A and B.

% ============================================================================
\section{The three jobs (A, B, C)}
% ============================================================================

\begin{table}[h!]
\centering
\renewcommand{\arraystretch}{1.35}
\begin{tabular}{@{}L{1.4cm}L{9.0cm}L{2.7cm}@{}}
\toprule
\textbf{Job} & \textbf{What it does} & \textbf{Status} \\
\midrule
\textbf{Kernel A} &
Trains the Boston detector on the full data and measures the untouched pipeline
once. This is the setup and timing check. &
\textbf{Done} \\
\textbf{Kernel B} &
Runs the main experiment: $15$ retraining runs of the detector (varying data
size, training length, and seed) and measures the effect on prediction and
planning. This produces the headline result. &
\textbf{Done} \\
\textbf{Kernel C} &
Repeats the $15$-style experiment with a class-balanced training set, to check
whether balancing the rare classes still gives the same cases. &
\textbf{Done ($7$ of $12$)} \\
\bottomrule
\end{tabular}
\end{table}

\noindent
The numbers in this report come from all three jobs. Kernel~C finished $7$ of its
$12$ planned runs; the other $5$ were the heaviest balancing setting and ran out
of cloud time, so I stopped there because the $7$ finished runs already answer the
question. Its results are in Section~\ref{sec:balanced}.

% ============================================================================
\section{The dataset}
% ============================================================================

The experiment now uses the full \textbf{nuScenes v1.0 trainval} dataset, the
standard public benchmark for this kind of work. It contains $850$ scenes:
$467$ recorded in Boston and $383$ in Singapore. We train the detector on Boston
and fine-tune it on Singapore, so the retraining is a real geographic domain
shift. The dataset (about $45$~GB) was already hosted on the cloud platform, so
no download or extra compute cost was needed.

The earlier version used only nuScenes-mini, a tiny $10$-scene demo subset. That
is the core difference: we moved from the demo subset to the full benchmark.

% ============================================================================
\section{The one change that drives everything}
% ============================================================================

The previous version of the experiment used a very small test set: only $3$
scenes, which is about $121$ camera frames. The new version uses a large test
set: $50$ scenes, which is $2{,}020$ frames. A scene is a $20$-second driving
clip and holds roughly $40$ frames, so the frame count is the fair way to
compare the two.

One point to be clear about: we train on all the Singapore data ($229$ scenes),
but we evaluate on $50$ of the $154$ Singapore validation scenes, chosen evenly
across the set. This keeps the running time manageable across $15$ retraining
runs while still giving a test set that is many times larger than before. The
$50$ scenes are fixed, so every run is measured on exactly the same frames.

Nothing about the method changed. Only the amount of data changed. Every result
below follows from that.

% ============================================================================
\section{How much data, before and after}
% ============================================================================

\begin{table}[h!]
\centering
\renewcommand{\arraystretch}{1.3}
\begin{tabular}{@{}L{4.6cm}C{3.7cm}C{4.6cm}@{}}
\toprule
 & \textbf{Before (mini)} & \textbf{After (full trainval)} \\
\midrule
Boston training scenes   & 6 scenes    & \textbf{280 scenes, 11{,}087 images} \\
Singapore fine-tune data & 119 images  & \textbf{229 scenes} \\
\textbf{Test set}        & \textbf{3 scenes, 121 frames} & \textbf{50 scenes, 2{,}020 frames} \\
\bottomrule
\end{tabular}
\end{table}

\noindent
The test set grew from $121$ frames to $2{,}020$ frames, roughly $17$ times
larger.

% ============================================================================
\section{The pipeline (unchanged)}
% ============================================================================

The system is the same three-part self-driving pipeline:

\begin{itemize}[noitemsep]
    \item \textbf{C1, Perception:} a camera detector (YOLOv11) that finds cars,
          pedestrians, and other agents.
    \item \textbf{C2, Prediction:} predicts where those agents will move.
    \item \textbf{C3, Planning:} decides the car's driving path.
\end{itemize}

We retrain only C1 on new-city data and measure what happens to C2 and C3. That
part did not change.

% ============================================================================
\section{The real numbers, side by side}
% ============================================================================

Both tables below report the same reference update: the detector is trained on
Boston, then fine-tuned on Singapore, and every component is measured on the
Singapore test set. ``Isolated'' means the component is fed perfect
ground-truth inputs; ``pipeline'' means it is fed the live output of the part
above it. The gap between the two is the cascade.

The last column reads each row as \textbf{better} (green sense: score moved the
good way) or \textbf{WORSE} (bad way). Entangled enhancement is the pattern where
C1 is better, C2 is better, and yet C3 is WORSE at the same time.

\subsection*{Before: mini test set (3 scenes, 121 frames)}

\begin{table}[h!]
\centering
\renewcommand{\arraystretch}{1.3}
\setlength{\tabcolsep}{6pt}
\begin{tabular}{@{}L{4.6cm}rrrl@{}}
\toprule
\textbf{Part / metric} & \textbf{Before} & \textbf{After} & \textbf{$\Delta$\%} & \textbf{Direction} \\
\midrule
\textbf{C1} detector mAP@50   & 0.1185  & 0.0644 & $-45.7$ & \textbf{WORSE} \\
\textbf{C2} pred., isolated   & 2.09    & 2.09   & $0.0$   & (frozen) \\
\textbf{C2} pred., pipeline   & 15.19   & 3.95   & $-74.0$ & better \\
\textbf{C3} plan, isolated    & 4.93    & 4.93   & $0.0$   & (frozen) \\
\textbf{C3} plan, pipeline    & 6.30    & 6.58   & $+4.5$  & \textbf{WORSE} \\
\textbf{C3} collision rate    & 0.00    & 0.00   & $0.0$   & none \\
\bottomrule
\end{tabular}
\end{table}

\noindent
Pattern here: \textbf{C1 WORSE}, C2 better, \textbf{C3 WORSE}. Because the
detector itself got worse, this is \emph{not} entangled enhancement (that needs
C1 to get better). This is why the mini paper had to add a class-balancing step
to force C1 up.

\subsection*{After: full test set (50 scenes, 2,020 frames)}

\begin{table}[h!]
\centering
\renewcommand{\arraystretch}{1.3}
\setlength{\tabcolsep}{6pt}
\begin{tabular}{@{}L{4.6cm}rrrl@{}}
\toprule
\textbf{Part / metric} & \textbf{Before} & \textbf{After} & \textbf{$\Delta$\%} & \textbf{Direction} \\
\midrule
\textbf{C1} detector mAP@50   & 0.137   & 0.256  & $+86.6$ & \textbf{better} \\
\textbf{C2} pred., isolated   & 1.12    & 1.12   & $0.0$   & (frozen) \\
\textbf{C2} pred., pipeline   & 7.93    & 5.76   & $-27.4$ & \textbf{better} \\
\textbf{C3} plan, isolated    & 4.64    & 4.64   & $0.0$   & (frozen) \\
\textbf{C3} plan, pipeline    & 5.55    & 5.67   & $+2.2$  & \textbf{WORSE} \\
\textbf{C3} collision rate    & 0.02    & 0.00   & $-100$  & \textbf{better} \\
\bottomrule
\end{tabular}
\end{table}

\noindent
Pattern here: \textbf{C1 better, C2 better, and the C3 planning error (L2) WORSE}.
This is exactly \textbf{entangled enhancement}, and it appears with no trick. The
small-data drop in C1 from the mini set is gone.

\vspace{0.4em}
\noindent
One important correction from the earlier reading. The two planning numbers do not
agree. On the full test set the collision rate actually \textbf{went down}, from
$2\%$ to $0\%$, at the same time as the L2 planning error went \textbf{up}. So the
detector got better, and the system got safer on collisions, but its path matched
the human driver a little less. The two system-level measures move in opposite
directions. This is why we do not say the update simply harms the system. We say
the terminal metrics disagree, which is the real reason a single number at the
output is not enough to accept or reject an update, and why the proposed method
reports the whole set of downstream numbers instead of one.

% ============================================================================
\section{The full result: all 15 retraining runs (after)}
% ============================================================================

This is the complete Kernel~B result. Each row is one retraining of the detector,
under a different setting. The first three columns are the settings we varied:
\textbf{Data} is how much Singapore data we fine-tuned on (full, half, or a
quarter), \textbf{Ep.} is how many training epochs, and \textbf{Seed} is the
random seed (repeating each setting rules out luck).

The three ``$\Delta$'' columns show whether each part improved or worsened
compared to the baseline: for C1 a \textbf{positive} $\Delta$ is better (higher
detection score); for C2 and C3 a \textbf{negative} $\Delta$ is better (lower
error). The last column marks the row \textbf{YES} when C1 improved and C3
worsened at the same time, which is entangled enhancement. Those rows are in
bold.

Baseline (before any retraining): C1 mAP $=0.137$, C2 pipeline $=7.93$~m,
C3 pipeline $=5.55$~m.

\begin{table}[h!]
\centering
\renewcommand{\arraystretch}{1.2}
\setlength{\tabcolsep}{4.5pt}
\small
\begin{tabular}{@{}lccrrrc@{}}
\toprule
\textbf{Data} & \textbf{Ep.} & \textbf{Seed}
& \textbf{C1 $\Delta$} & \textbf{C2 $\Delta$ (m)} & \textbf{C3 $\Delta$ (m)}
& \textbf{Ent.\ enh.?} \\
 & & & \footnotesize (+ better) & \footnotesize ($-$ better) & \footnotesize ($-$ better) & \\
\midrule
\textbf{Full} & \textbf{10} & \textbf{1} & \textbf{$+0.119$} & \textbf{$-2.17$} & \textbf{$+0.124$} & \textbf{YES} \\
\textbf{Full} & \textbf{10} & \textbf{2} & \textbf{$+0.122$} & \textbf{$-2.28$} & \textbf{$+0.193$} & \textbf{YES} \\
\textbf{Full} & \textbf{10} & \textbf{3} & \textbf{$+0.129$} & \textbf{$-2.15$} & \textbf{$+0.187$} & \textbf{YES} \\
Half & 10 & 1 & $+0.114$ & $-2.55$ & $-0.121$ & no \\
Half & 10 & 2 & $+0.120$ & $-2.71$ & $-0.230$ & no \\
Half & 10 & 3 & $+0.106$ & $-1.77$ & $-0.151$ & no \\
\textbf{Qtr} & \textbf{10} & \textbf{1} & \textbf{$+0.121$} & \textbf{$-2.71$} & \textbf{$+0.073$} & \textbf{YES} \\
Qtr & 10 & 2 & $+0.113$ & $-1.29$ & $-0.093$ & no \\
\textbf{Qtr} & \textbf{10} & \textbf{3} & \textbf{$+0.101$} & \textbf{$-2.82$} & \textbf{$+0.008$} & \textbf{YES} \\
Full & 5 & 1 & $+0.138$ & $-1.08$ & $-0.095$ & no \\
Full & 5 & 2 & $+0.122$ & $-1.88$ & $-0.130$ & no \\
Full & 5 & 3 & $+0.130$ & $-1.08$ & $-0.041$ & no \\
\textbf{Full} & \textbf{20} & \textbf{1} & \textbf{$+0.142$} & \textbf{$-2.59$} & \textbf{$+0.037$} & \textbf{YES} \\
Full & 20 & 2 & $+0.126$ & $-1.77$ & $-0.083$ & no \\
Full & 20 & 3 & $+0.130$ & $-2.44$ & $-0.020$ & no \\
\midrule
\multicolumn{3}{@{}l}{\textbf{Totals}} & \textbf{15/15 better} & \textbf{15/15 better} & \textbf{6/15 worse} & \textbf{6 YES} \\
\bottomrule
\end{tabular}
\end{table}

\noindent
How to read it in one look: the C1 column is positive in every row (detector
always improved), the C2 column is negative in every row (prediction always
improved), and the C3 column is positive in $6$ rows (planner worse). Those $6$
rows are the entangled-enhancement cases: a better detector, a better predictor,
and yet a worse driving plan.

% ============================================================================
\section{The class-balanced check (Kernel C)}
\label{sec:balanced}
% ============================================================================

In the old small-scale experiment, we had to balance the rare classes (oversample
the images that contain trucks, buses, and so on) just to make the detector score
go up at all. On the full data that step is no longer needed, because Kernel~B
already shows the effect. So here Kernel~C asks a different question: if we still
balance the classes, does the entangled enhancement survive, or was it an artifact
of the class imbalance?

I ran two balancing strengths, two training lengths, and three seeds ($12$ planned
runs). Seven finished; the other five were the heaviest setting and ran out of
cloud time. The seven finished runs are enough to answer the question.

\begin{table}[h!]
\centering
\renewcommand{\arraystretch}{1.2}
\setlength{\tabcolsep}{5pt}
\small
\begin{tabular}{@{}lccrrc@{}}
\toprule
\textbf{Balance} & \textbf{Ep.} & \textbf{Seed}
& \textbf{C1 $\Delta$} & \textbf{C3 $\Delta$ (m)} & \textbf{Ent.\ enh.?} \\
 & & & \footnotesize (+ better) & \footnotesize ($-$ better) & \\
\midrule
Light ($p{=}0.5$) & 10 & 1 & $+0.127$ & $-0.021$ & no \\
Light ($p{=}0.5$) & 10 & 2 & $+0.122$ & $-0.023$ & no \\
\textbf{Light ($p{=}0.5$)} & \textbf{10} & \textbf{3} & \textbf{$+0.112$} & \textbf{$+0.014$} & \textbf{YES} \\
\textbf{Light ($p{=}0.5$)} & \textbf{20} & \textbf{1} & \textbf{$+0.141$} & \textbf{$+0.139$} & \textbf{YES} \\
\textbf{Light ($p{=}0.5$)} & \textbf{20} & \textbf{2} & \textbf{$+0.143$} & \textbf{$+0.094$} & \textbf{YES} \\
Light ($p{=}0.5$) & 20 & 3 & $+0.137$ & $-0.001$ & no \\
\textbf{Heavy ($p{=}1.0$)} & \textbf{20} & \textbf{1} & \textbf{$+0.125$} & \textbf{$+0.075$} & \textbf{YES} \\
\midrule
\multicolumn{3}{@{}l}{\textbf{Totals}} & \textbf{7/7 better} & \textbf{4/7 worse} & \textbf{4 YES} \\
\bottomrule
\end{tabular}
\end{table}

\noindent
The result: even with class balancing, the detector still improves in every run
($7$ of $7$), and entangled enhancement still appears ($4$ of $7$). So the effect
is not caused by class imbalance. It is a real property of the pipeline under
retraining, and balancing the classes does not remove it.

% ============================================================================
\section{Re-checking CARA's cheap early-warning screen}
% ============================================================================

The proposed method has a cheap first-pass screen: instead of running the whole
pipeline, it looks only at how much the new detector's outputs drifted from the old
one, and uses that drift score to guess which updates are risky. On the small test
set this screen looked promising. On the full test set I re-checked it properly,
and the honest answer is more limited.

\begin{table}[h!]
\centering
\renewcommand{\arraystretch}{1.3}
\begin{tabular}{@{}L{7.2cm}C{3.0cm}C{3.0cm}@{}}
\toprule
\textbf{What we asked of the drift score} & \textbf{Before (mini)} & \textbf{After (full)} \\
\midrule
Does it track \emph{how much} the plan moves? (rank correlation) & about $0.41$ & \textbf{about $0.56$} \\
Does it tell a worse update from a safe one? (AUROC) & about $0.70$ & \textbf{about $0.48$ (chance)} \\
\bottomrule
\end{tabular}
\end{table}

\noindent
So the drift screen does track the \emph{size} of the disturbance, and a little
better than before ($0.56$). But it cannot tell the \emph{direction}: it does not
separate the updates that hurt the planning metric from the ones that help. On the
full data all the drift scores fall in a narrow band, so the risky and safe updates
sit mixed together. The earlier $0.70$ was small-sample optimism.

I also validated it the careful way this time, choosing the threshold on one half of
the updates and testing on the untouched other half ($8$ updates). On that held-out
test it still catches every planning-metric regression, but with low precision, and
with only $8$ updates the error bars are wide (for example the recall could sit
anywhere from $0.44$ to $1.00$). The takeaway is honest and useful: the cheap screen
is a good way to \emph{prioritize} which updates to check first, but it is not a
stand-alone accept or reject rule. The full check that runs the whole pipeline
remains the reliable one.

% ============================================================================
\section{Summary of what changed}
% ============================================================================

\begin{table}[h!]
\centering
\renewcommand{\arraystretch}{1.3}
\begin{tabular}{@{}L{6.3cm}C{3.2cm}C{3.5cm}@{}}
\toprule
\textbf{What we measured} & \textbf{Before (3 sc., 121 fr.)} & \textbf{After (50 sc., 2{,}020 fr.)} \\
\midrule
Detector improves after retraining?           & 0 of 15 & \textbf{15 of 15} \\
Prediction improves?                           & mixed   & \textbf{15 of 15} \\
\textbf{Entangled enhancement seen directly?}  & \textbf{0 of 15 (needed a trick)} & \textbf{6 of 15 (no trick)} \\
Baseline collision rate                        & 0\% & \textbf{2\%} \\
Collision rate after retraining                & 0\% & \textbf{0\% (fell from 2\%)} \\
Plan shift range                               & 1.33 to 7.80 m & 0.78 to 1.50 m \\
\bottomrule
\end{tabular}
\end{table}

\noindent
The plan-shift range also became tighter, from $1.33$--$7.80$~m down to
$0.78$--$1.50$~m. This is expected: with a much larger test set the measurement
is less noisy, so the extreme values from the tiny $3$-scene set do not appear.

% ============================================================================
\section{Baseline system, before any retraining}
% ============================================================================

For reference, the untouched pipeline measured on the $2{,}020$-frame test set:

\begin{table}[h!]
\centering
\renewcommand{\arraystretch}{1.3}
\begin{tabular}{@{}L{5.5cm}C{2.2cm}L{5.7cm}@{}}
\toprule
\textbf{Metric} & \textbf{Value} & \textbf{Reading} \\
\midrule
C1 detector mAP                 & 0.137 & detector works \\
C2 prediction, isolated         & 1.12 m & clean prediction \\
C2 prediction, pipeline         & 7.93 m & error inherited from the detector \\
C3 planning, isolated           & 4.64 m & planner alone \\
C3 planning, pipeline           & 5.55 m & \textbf{worse than isolated} \\
C3 collision rate               & 2\% & collisions now occur \\
\bottomrule
\end{tabular}
\end{table}

\noindent
The planner error in the pipeline ($5.55$ m) is larger than the planner error in
isolation ($4.64$ m). That gap is the cascade: the planner inherits the error
that flows down from perception. It now holds across $50$ scenes instead of $3$.

% ============================================================================
\section{What this means for the paper}
% ============================================================================

\begin{enumerate}[leftmargin=*]
    \item The main claim is stronger. Entangled enhancement in a real driving
          pipeline is now shown by $6$ ordinary retraining runs on full data,
          rather than by a special class-balancing step on $3$ scenes.

    \item One old finding is removed. The ``detector score drops after
          retraining'' result was caused by the small $119$-image fine-tune set.
          With full data it does not happen, so that finding is dropped.

    \item A new and more careful result appears. At full scale the baseline
          pipeline has a $2\%$ collision rate (there were none on $3$ scenes), and
          retraining drives it to zero while the L2 planning error goes up. The two
          system-level numbers move in opposite directions, so the finding is
          stated as a disagreement between metrics, not as a plain harm. This makes
          the point of the proposed method clearer: it must report the whole set of
          downstream numbers, because no single one is safe to trust alone.

    \item The cheap early-warning screen was re-checked honestly. It tracks the
          size of the disturbance but not its direction, so it is a way to
          prioritize which updates to check, not a stand-alone accept or reject
          rule.
\end{enumerate}

% ============================================================================
\section{What is left}
% ============================================================================

All three cloud jobs are finished, so the experiment itself is complete and every
number in this report is final. Two small things remain on the writing side, and
neither needs new experiments to state:

\begin{itemize}[leftmargin=*]
    \item The five class-balanced runs that ran out of cloud time were the heaviest
          setting. The seven finished runs already show the effect survives
          balancing, so finishing the last five is optional and would only add more
          of the same.
    \item One check would strengthen the smallest of the six entangled-enhancement
          cases (the one where the planning error rose by only $0.008$~m): running
          the unchanged detector against itself to measure the natural measurement
          noise. If that noise is below $0.008$~m, all six cases stand; if not, that
          one case is borderline. This is a quick run I can do next.
\end{itemize}

\end{document}
